\section{Contribuciones y límites de los LLM: el espacio virtual no es todo el mundo}

Ya se estableció la definición de modelo del mundo; ahora hay que distinguirlo de los LLM. Esta distinción es importante, porque la entrevista no está en contra de los LLM, sino que pregunta qué parte del mundo cubre realmente el éxito de los LLM.

Xie Saining (谢赛宁) no niega el carácter revolucionario de los LLM. Dice que tiene que agradecer a los LLM, porque sin ellos la inteligencia multimodal no habría escalado hasta su tamaño actual. Pero se opone a extender el relato de los LLM hasta afirmar que «los modelos de lenguaje conducen de manera natural a la human-level intelligence». En su marco, los LLM son más fuertes en el digital / virtual space: texto, código, síntesis de conocimiento, educación, derecho, búsqueda, agentic coding, etc. Pueden ser un elemento importante de un sistema inteligente, pero no son el fundamento del modelo del mundo.

\begin{figure}[H]
\centering
\includegraphics[width=0.96\textwidth]{llm-vs-world-model.png}
\caption{LLM y modelos del mundo: comparación de dos mercados de la inteligencia.}
\end{figure}

\begin{knowledgebox}{Lectura de la figura: los modelos de lenguaje y los modelos del mundo se complementan, pero no se sustituyen entre sí}
A la izquierda, los puntos fuertes del LLM son el conocimiento tokenizado, el razonamiento sobre texto y la operación en el espacio digital; a la derecha, el objetivo del world model son las señales continuas, la predicción física, la planificación de acciones y el control seguro. Ambos pueden complementarse: el lenguaje aporta la interfaz, el conocimiento y la comunicación; el modelo del mundo aporta el anclaje físico, la predicción dinámica y la capacidad de actuar. Reducir a la fuerza uno al otro lleva a diagnosticar mal los cuellos de botella tecnológicos.
\end{knowledgebox}

\subsection{Por qué los modelos de lenguaje se parecen al aprendizaje con supervisión fuerte}

En la entrevista aparece una idea muy reveladora: a los modelos de lenguaje se les suele describir como self-supervised learning, pero desde otro ángulo, el lenguaje mismo ya es una señal de supervisión fuerte, procesada por la civilización humana. Miles de años de civilización, libros, páginas web, código, artículos científicos e internet han comprimido una enorme cantidad de conocimiento del mundo en texto tokenizado. Entrenar un LLM se parece a descargar ese conocimiento ya procesado, no a aprender directamente del mundo físico.

Esto explica por qué la scaling law de los LLM pudo aparecer relativamente pronto: cuenta con una gran cantidad de material que los seres humanos ya organizaron, etiquetaron y volvieron comunicable. En cambio, vision y robotics trabajan con raw sensory data, acciones, espacio, restricciones físicas y retroalimentación, y por naturaleza carecen de «etiquetas al estilo de internet» de la misma escala, la misma calidad y el mismo grado de compresión.

\begin{warningbox}{«Datos gratuitos» no equivale a «datos sin supervisión»}
Que el texto de internet pueda recopilarse gratis no significa que carezca de etiquetas humanas. El lenguaje es una estructura que los seres humanos comprimieron durante largo tiempo para comunicarse; cuando el modelo ingiere ese texto, también ingiere la gran cantidad de abstracción, selección y etiquetado que los seres humanos ya realizaron.
\end{warningbox}

\subsection{Una taza que se rompe: lo que omite la descripción lingüística}

Xie Saining da un ejemplo sencillo: cuando decimos «la taza cayó al suelo y se rompió», el lenguaje solo conserva la información necesaria para comunicarse. No describe el contacto de la taza, las fuerzas que actúan sobre ella, la trayectoria de la fractura, las propiedades del material, la dispersión de los fragmentos ni el sonido. Para la comunicación humana, la mayoría de esos detalles son innecesarios; para un sistema que debe actuar en el mundo físico, algunos de ellos pueden ser cruciales.

Por lo tanto, el lenguaje no es el mundo mismo, sino una compresión hecha para comunicarse. Lo que el modelo del mundo necesita aprender son precisamente las regularidades que quedan fuera de esa compresión lingüística: dinámica, estructura espacial, contacto, secuencia temporal, causalidad y posibilidad de acción.

\subsection{Resumen de la sección}

El éxito de los LLM proviene de la enorme compresión del conocimiento humano y de la escala de internet. Seguirán siendo un componente importante de los sistemas inteligentes, pero la inteligencia orientada al mundo físico no puede depender solo del lenguaje. Lo que el modelo del mundo debe completar es la parte del mundo que el lenguaje omite por naturaleza.
